\documentclass[11pt,letterpaper]{article}

\usepackage[top=1in, bottom=1in, left=1in, right=1in, headheight=14pt]{geometry}
\usepackage{fontspec}
\setmainfont{DejaVu Serif}
\setsansfont{DejaVu Sans}
\setmonofont{DejaVu Sans Mono}

\usepackage{xcolor}
\usepackage{titlesec}
\usepackage{fancyhdr}
\usepackage{enumitem}
\usepackage{hyperref}
\usepackage{booktabs}
\usepackage{tabularx}
\usepackage{longtable}
\usepackage{colortbl}
\usepackage{multirow}
\usepackage{array}
\usepackage{parskip}
\usepackage{tcolorbox}
\usepackage{graphicx}
\usepackage{lastpage}
\usepackage{microtype}

% ── Color Scheme: Puget Sound / Bioacoustics ─────────────────────────────────
\definecolor{primary}{HTML}{004D40}       % Deep teal — Puget Sound depth
\definecolor{secondary}{HTML}{1A237E}     % Deep navy — night forest / Part A echo
\definecolor{tertiary}{HTML}{33691E}      % Fern green — old-growth understory
\definecolor{accent}{HTML}{0277BD}        % Sky blue — Salish Sea surface
\definecolor{partE}{HTML}{4A148C}         % Deep purple — synthesis/completion
\definecolor{lightprimary}{HTML}{E0F2F1}
\definecolor{lightsecondary}{HTML}{E8EAF6}
\definecolor{lighttertiary}{HTML}{F1F8E9}
\definecolor{lightE}{HTML}{F3E5F5}
\definecolor{rowgray}{HTML}{F5F5F5}
\definecolor{bordergray}{HTML}{BDBDBD}
\definecolor{subtlegray}{HTML}{757575}
\definecolor{darktext}{HTML}{212121}
\definecolor{blockred}{HTML}{B71C1C}
\definecolor{gold}{HTML}{F57F17}
\definecolor{coral}{HTML}{BF360C}

% ── Color aliases for each part ───────────────────────────────────────────────
\definecolor{colorA}{HTML}{1A237E}   % Part A — navy (music)
\definecolor{colorB}{HTML}{004D40}   % Part B — teal (nature)
\definecolor{colorC}{HTML}{E65100}   % Part C — amber (synthesis/fire)
\definecolor{colorD}{HTML}{4527A0}   % Part D — violet (fractal depth)
\definecolor{colorE}{HTML}{1B5E20}   % Part E — dark green (completion/growth)

% ── Typography ─────────────────────────────────────────────────────────────────
\titleformat{\section}
  {\Large\bfseries\sffamily\color{primary}}
  {\thesection}{1em}{}
  [\vspace{-0.5em}\textcolor{bordergray}{\rule{\textwidth}{0.4pt}}]

\titleformat{\subsection}
  {\large\bfseries\sffamily\color{secondary}}
  {\thesubsection}{1em}{}

\titleformat{\subsubsection}
  {\normalsize\bfseries\sffamily\color{tertiary}}
  {\thesubsubsection}{1em}{}

\titlespacing*{\section}{0pt}{1.5em}{0.8em}
\titlespacing*{\subsection}{0pt}{1.2em}{0.5em}
\titlespacing*{\subsubsection}{0pt}{0.8em}{0.3em}

% ── Custom Environments ────────────────────────────────────────────────────────
\tcbuselibrary{breakable,skins}

\newcommand{\stageheader}[2]{%
  \noindent\colorbox{#2}{\makebox[\textwidth][l]{%
    \textcolor{white}{\bfseries\sffamily\large\hspace{8pt}#1\hspace{8pt}}}}%
  \vspace{0.6em}\par%
}

\newcommand{\partheader}[3]{%
  \noindent\colorbox{#3}{\makebox[\textwidth][l]{%
    \textcolor{white}{\bfseries\sffamily\Large\hspace{10pt}PART #1 \quad #2\hspace{10pt}}}}%
  \vspace{0.8em}\par%
}

\newtcolorbox{asciibox}{
  colback=rowgray, colframe=bordergray,
  arc=1pt, boxrule=0.5pt,
  left=6pt, right=6pt, top=4pt, bottom=4pt,
  fontupper=\small\ttfamily,
  breakable
}

\newtcolorbox{quotebox}{
  colback=lightprimary, colframe=primary,
  arc=2pt, boxrule=0.8pt,
  left=12pt, right=12pt, top=8pt, bottom=8pt,
  breakable
}

\newtcolorbox{warnbox}{
  colback=lightsecondary, colframe=secondary,
  arc=2pt, boxrule=0.8pt,
  left=12pt, right=12pt, top=8pt, bottom=8pt,
  breakable
}

\newtcolorbox{greenbox}{
  colback=lighttertiary, colframe=tertiary,
  arc=2pt, boxrule=0.8pt,
  left=12pt, right=12pt, top=8pt, bottom=8pt,
  breakable
}

\newtcolorbox{amberbox}{
  colback=lightsecondary, colframe=colorC,
  arc=2pt, boxrule=0.8pt,
  left=12pt, right=12pt, top=8pt, bottom=8pt,
  breakable
}

\newcommand{\metafield}[2]{\textbf{\sffamily\textcolor{primary}{#1}} \textcolor{darktext}{#2}\par}
\newcolumntype{L}[1]{>{\raggedright\arraybackslash}p{#1}}
\newcolumntype{C}[1]{>{\centering\arraybackslash}p{#1}}
\newcommand{\tblhdr}[1]{\textbf{\sffamily\textcolor{white}{#1}}}

% ── Headers / Footers ─────────────────────────────────────────────────────────
\pagestyle{fancy}
\fancyhf{}
\fancyhead[L]{\small\sffamily\textcolor{subtlegray}{The Sound of Puget Sound — A·B·C·D·E}}
\fancyhead[R]{\small\sffamily\textcolor{subtlegray}{GSD Mission Package}}
\fancyfoot[C]{\small\sffamily\textcolor{subtlegray}{Page \thepage\ of \pageref{LastPage}}}
\renewcommand{\headrulewidth}{0.4pt}
\renewcommand{\footrulewidth}{0pt}

\hypersetup{
  colorlinks=true, linkcolor=accent, urlcolor=accent,
  pdfauthor={GSD Ecosystem / Miles Tiberius Foxglove},
  pdftitle={The Sound of Puget Sound -- Complete A-B-C-D-E Mission Package},
}

% ═══════════════════════════════════════════════════════════════════════════════
\begin{document}

% ── Title Page ─────────────────────────────────────────────────────────────────
\thispagestyle{empty}
\begin{center}
\vspace*{1.5cm}

{\small\sffamily\textcolor{primary}{\textbf{GSD RESEARCH MISSION PACKAGE · FLEET+ ACTIVATION}}}

\vspace{0.3cm}
\textcolor{bordergray}{\rule{10cm}{0.4pt}}

\vspace{1cm}
{\fontsize{34}{40}\selectfont\bfseries\sffamily\textcolor{primary}{The Sound of Puget Sound}}

\vspace{0.5cm}
{\fontsize{14}{18}\selectfont\sffamily\textcolor{secondary}{Complete A–B–C–D–E Release Matrix Mission Package}}

\vspace{0.8cm}
{\large\sffamily\textcolor{tertiary}{360 Musicians · 360 Wild Voices · One Unit Circle}}

\vspace{0.3cm}
{\normalsize\sffamily\itshape\textcolor{subtlegray}{From Quincy Jones to Puget Sound in Silence · Degree by Degree}}

\vspace{1.5cm}
% Part labels
\begin{tabular}{ccccc}
\colorbox{colorA}{\textcolor{white}{\sffamily\bfseries\small\  PART A \ }} &
\colorbox{colorB}{\textcolor{white}{\sffamily\bfseries\small\  PART B \ }} &
\colorbox{colorC}{\textcolor{white}{\sffamily\bfseries\small\  PART C \ }} &
\colorbox{colorD}{\textcolor{white}{\sffamily\bfseries\small\  PART D \ }} &
\colorbox{colorE}{\textcolor{white}{\sffamily\bfseries\small\  PART E \ }} \\[4pt]
\sffamily\tiny\textcolor{subtlegray}{Seattle 360} &
\sffamily\tiny\textcolor{subtlegray}{PNW Species} &
\sffamily\tiny\textcolor{subtlegray}{Deep Mapping} &
\sffamily\tiny\textcolor{subtlegray}{Fractal Refine} &
\sffamily\tiny\textcolor{subtlegray}{Lessons Applied} \\
\end{tabular}

\vspace{1.5cm}
{\sffamily\textcolor{accent}{Vision $\rightarrow$ Research $\rightarrow$ Mission $\rightarrow$ Refinement $\rightarrow$ Synthesis}}

\vspace{1.5cm}
\begin{tabular}{ll}
\textcolor{subtlegray}{\sffamily\small Date:} & \textcolor{darktext}{\sffamily\small 2026-03-28} \\[3pt]
\textcolor{subtlegray}{\sffamily\small Inputs:} & \textcolor{darktext}{\sffamily\small seattle\_360\_geo.csv (360 artists) + pnw\_360\_species.csv (360 species)} \\[3pt]
\textcolor{subtlegray}{\sffamily\small Total Releases:} & \textcolor{darktext}{\sffamily\small 1,800 artifacts across Parts A–E} \\[3pt]
\textcolor{subtlegray}{\sffamily\small Components:} & \textcolor{darktext}{\sffamily\small 37 components (9+9+7+6+6) across 5 parts} \\[3pt]
\textcolor{subtlegray}{\sffamily\small Total Tests:} & \textcolor{darktext}{\sffamily\small 570 (148+152+120+90+60)} \\[3pt]
\textcolor{subtlegray}{\sffamily\small Token Budget:} & \textcolor{darktext}{\sffamily\small $\approx$77.6M total (A: 20.7M · B: 22.3M · C: 13.3M · D: 14.4M · E: 4.8M)} \\[3pt]
\textcolor{subtlegray}{\sffamily\small Fleet Activation:} & \textcolor{darktext}{\sffamily\small Full crew · PNW species naming · OCAP®/CARE/UNDRIP compliant} \\[3pt]
\textcolor{subtlegray}{\sffamily\small College Depts:} & \textcolor{darktext}{\sffamily\small 15+ departments seeded including new: science/bioacoustics/, diy-art/} \\[3pt]
\textcolor{subtlegray}{\sffamily\small Status:} & \textcolor{darktext}{\sffamily\small Ready for GSD Orchestrator} \\
\end{tabular}

\vspace{1.5cm}
\textcolor{bordergray}{\rule{9cm}{0.4pt}}\\[0.3em]
{\small\sffamily\textcolor{subtlegray}{GSD Ecosystem · github.com/Tibsfox · Miles Tiberius Foxglove · 2026}}
\end{center}
\newpage

\tableofcontents
\newpage

% ─────────────────────────────────────────────────────────────────────────────
% EXECUTIVE OVERVIEW
% ─────────────────────────────────────────────────────────────────────────────
\section{Executive Overview}

\begin{quotebox}
\textit{``Before there were amplifiers, there was the Salish Sea. Before there were drum machines,
there was the Roosevelt elk bugling at dawn in the Hoh River valley. Before there was a music
theory textbook, there was the orca pod's greeting ceremony. The Pacific Northwest is not a
backdrop to human culture. It is a teacher.''}
\end{quotebox}

\vspace{0.8em}

The Sound of Puget Sound is a five-part autonomous educational engine that processes 360
Seattle musicians (Part A) and 360 PNW species and soundscapes (Part B) through the same
Fleet research pipeline — both mapped to the same 360° unit circle — then weaves their
outputs into a unified College of Knowledge curriculum (Part C), refines every release to
publication quality (Part D), and synthesizes all lessons learned into the specification
for the next complete rotation (Part E).

\subsection{The Parallel Structure}

Every degree on the circle has two entries: a human musical voice and a wild biological voice.
They are not metaphors for each other. They are structural parallels — exhibiting the same
acoustic, mathematical, and cultural patterns in different substrates.

\begin{center}
\rowcolors{2}{rowgray}{white}
\begin{tabularx}{\textwidth}{C{1.2cm}L{4cm}L{4.5cm}X}
\toprule
\rowcolor{primary}
\tblhdr{Degree} & \tblhdr{Part A (Music)} & \tblhdr{Part B (Species)} & \tblhdr{Shared Principle} \\
\midrule
0° & Quincy Jones (Jazz, energy 1–3) & Great Blue Heron (silence/croak, E=1) & Call-and-response; structured silence \\
1° & Bill Frisell (chromaticism) & Pacific Wren (micro-phrase cascade) & Complex melodic improvisation \\
45° & Damien Jurado (folk) & Pacific Chorus Frog (chorus) & Rhythmic entrainment; unison emergence \\
108° & Pileated Woodpecker (drum) & Pileated Woodpecker (same!) & This degree IS the woodpecker \\
276° & Bikini Kill (riot grrrl, E=8) & Orca female-led hunt (E=8) & Matriarchal coordination as structure \\
311° & Nirvana (grunge, E=9) & Orca greeting ceremony (E=9) & Full community convergence call \\
344° & Botch (mathcore, E=10) & Multi-species polyrhythm apex (E=10) & Prime-number meters; natural mathematics \\
359° & Unwound (silence, E=10) & Puget Sound one breath (E=10) & Silence as composition; the circle closes \\
\bottomrule
\end{tabularx}
\end{center}

\subsection{The Unit Circle as Structure}

\begin{asciibox}
                    0deg: Jazz (low energy) | Great Blue Heron
                           ^
         Marine/Coastal    |    Freshwater/Riparian
         + Folk/Soul       |    + Indie/Electronic
         (0-89deg)         |    (90-179deg)
                           |
  270deg <-----------------+-------------------> 90deg
  Grunge/Metal             |             Folk/Singer-Songwriter
  + Storm/Geological       |             + Forest Species
  (270-359deg)             |             (90-179deg)
                           |
                           v
              180deg: Blues/Soul | Salmon + Owls + Mammals

  Inner ring (energy 1-3): quiet, foundational, subtle
  Middle ring (energy 4-6): intermediate, complex, layered  
  Outer ring (energy 7-10): intense, advanced, apex
  
  BOTH DATASETS use the same circle. Same mathematics. Different language.
\end{asciibox}

\subsection{Five-Part Summary}

\begin{center}
\rowcolors{2}{rowgray}{white}
\begin{tabularx}{\textwidth}{C{1cm}L{2.5cm}L{3cm}C{2cm}C{2cm}X}
\toprule
\rowcolor{primary}
\tblhdr{Part} & \tblhdr{Name} & \tblhdr{What It Produces} & \tblhdr{Releases} & \tblhdr{Tokens} & \tblhdr{Status} \\
\midrule
A & Seattle 360 & 360 artist research packs & 360 & $\sim$20.7M & Delivered \\
B & PNW 360 Species & 360 bioacoustic packs & 360 & $\sim$22.3M & This package \\
C & Deep Mapping & 360 A×B synthesis docs + DIY art & 360 pairs & $\sim$13.3M & After A+B \\
D & Fractal Refine & Quality passes on 720 releases & 720 passes & $\sim$14.4M & After C \\
E & Lessons Applied & Meta-retrospective + Cycle 2 spec & 1 synthesis & $\sim$4.8M & After D \\
\midrule
\rowcolor{lightprimary}
& \textbf{Total} & \textbf{1,800 artifacts} & \textbf{1,800} & \textbf{$\sim$75.5M} & \\
\bottomrule
\end{tabularx}
\end{center}

\newpage

% ─────────────────────────────────────────────────────────────────────────────
% PART B — PNW 360 SPECIES
% ─────────────────────────────────────────────────────────────────────────────
\partheader{B}{PNW 360 Species Acoustic Release Engine}{colorB}

\section{Part B Vision: PNW Wild Voices}

\metafield{Input:}{\texttt{pnw\_360\_species.csv} — 360 species / soundscapes, degrees 0°–359°}
\metafield{Parallel to:}{Part A (same loop architecture, same College integration, same Safety Warden)}
\metafield{New departments:}{\texttt{.college/science/bioacoustics/} · \texttt{.college/diy-art/}}

\vspace{0.5em}

The orca pod's greeting ceremony has been formally analyzed by bioacousticians and found to
exhibit properties identical to human polyphonic choral music: distinct voice roles, shared
repertoire, call-and-response structure, and cultural transmission across generations.
The Pacific Wren sings two-voice polyphony simultaneously from a single syrinx. The stonefly
hatch at degree 329 produces a rhythmic wall of sound with polyrhythmic sub-patterns at the
prime-number ratios Botch would use sixty years later at degree 332.

Part B makes this explicit, sourced, and teachable. 360 species. 360 Fleet research missions.
One bioacoustic theory genealogy that traces sound from a heron's silence to the end of Puget Sound.

\subsection{PNW 360 Species CSV — Sample Entries}

\begin{center}
\rowcolors{2}{rowgray}{white}
\begin{tabularx}{\textwidth}{C{1.2cm}L{4.5cm}L{3.5cm}C{1.2cm}X}
\toprule
\rowcolor{colorB}
\tblhdr{Deg} & \tblhdr{Species / Soundscape} & \tblhdr{Sound Type} & \tblhdr{Energy} & \tblhdr{Ecosystem} \\
\midrule
0 & Great Blue Heron & stalking silence / pterodactyl croak & 1 & Puget Sound estuary \\
1 & Pacific Wren & intricate cascade warble & 2 & old-growth understory \\
45 & Pacific Chorus Frog & krek-ek breeding chorus & 2 & pond wetland \\
57 & Pacific Giant Salamander & low bark / growl when threatened & 3 & cold stream \\
100 & Common Raven & deep cronk + bell + hollow knock & 4 & open forest \\
108 & Pileated Woodpecker & wild laughing wuk-wuk-wuk & 6 & old-growth \\
116 & American Dipper & loud bubbly song over rushing water & 5 & fast stream \\
191 & Humpback Whale song & complex hour-long spiral song & 8 & open ocean \\
210 & Gray Wolf howl & sustained bass resonant howl & 7 & remote forest \\
212 & Cougar scream & blood-curdling rising shriek & 8 & forest \\
246 & Waterfall Snoqualmie & roaring white-noise wall & 7 & basalt cliff \\
270 & SRKW S01 pod call & repertoire-specific clan calls & 8 & Salish Sea \\
276 & Orca echolocation burst & rapid click trains ultra-precise & 8 & hunting \\
311 & Orca greeting ceremony & full pod multi-voice convergence & 9 & Salish Sea \\
332 & Woodpecker+River+Frog polyrhythm & three independent pulse streams & 9 & riparian \\
344 & Multi-species polyrhythm apex & woodpecker+orca+thunder+frog & 10 & convergence \\
359 & Puget Sound in silence & one breath before the next wave & 10 & full Salish Sea \\
\bottomrule
\end{tabularx}
\end{center}

\subsection{Acoustic Theory Mapping Framework}

The core innovation of Part B is the systematic mapping of bioacoustic properties to
music theory concepts. Every species receives an \texttt{AcousticTheoryNodeList} with:

\begin{center}
\rowcolors{2}{rowgray}{white}
\begin{tabularx}{\textwidth}{L{4cm}XL{4.5cm}}
\toprule
\rowcolor{colorB}
\tblhdr{Acoustic Feature} & \tblhdr{Music Theory Parallel} & \tblhdr{College Path} \\
\midrule
Fundamental frequency (Hz) & Pitch / note value & \texttt{.college/music/theory/harmony/} \\
Harmonic series (overtones) & Timbre / chord color & \texttt{.college/mathematics/calculus/fourier/} \\
Call duration (ms) & Note duration / rhythm & \texttt{.college/music/theory/rhythm/} \\
Inter-call interval & Rest / silence as composition & \texttt{.college/music/theory/form/} \\
Solo/duet/chorus & Monophony/counterpoint/polyphony & \texttt{.college/music/theory/counterpoint/} \\
Frequency modulation & Vibrato / portamento / glissando & \texttt{.college/music/theory/scales/} \\
Amplitude envelope & ADSR (synthesis) & \texttt{.college/music/technology/synthesis/} \\
Acoustic synchrony & Rhythmic entrainment / groove & \texttt{.college/music/theory/rhythm/groove.md} \\
Echolocation & Sampling theorem / Nyquist & \texttt{.college/mathematics/calculus/fourier/} \\
Infrasound ($<$20 Hz) & Sub-bass / felt not heard & \texttt{.college/music/technology/synthesis.md} \\
\bottomrule
\end{tabularx}
\end{center}

\subsection{Part B Component Architecture}

\begin{asciibox}
pnw-360-engine/
|
+-- B00-shared-types -------- SpeciesProfile, AcousticNode, SoundscapeLayer schemas
+-- B01-species-intake ------- pnw_360_species.csv -> SpeciesProfile[360]
|                              Flags: type=SPECIES vs SOUNDSCAPE_COMPOSITION
+-- B02-acoustic-mapper ------ AcousticTheoryNodeList (Wave 1A / EAGLE / Sonnet)
|                              Freq, rhythm, call structure, music parallel
+-- B03-fleet-research-gen --- Full LaTeX PDF per species (Wave 2 / HAWK / Opus)
|                              Color: deep teal / navy / fern green
+-- B04-college-linker ------- CollegeLinkList (Wave 1B / ORCA / Sonnet)
|                              New depts: science/bioacoustics/, diy-art/
+-- B05-release-pipeline ----- Sequential atomic publication (HERON / Sonnet)
+-- B06-retrospective -------- NASA SE lessons learned (OWL / Opus)
+-- B07-carry-forward -------- KnowledgeState management (CEDAR / Sonnet)
+-- B08-safety-warden -------- OCAP + ESA + location gates (WARDEN / Opus)
                               Extra gates vs Part A:
                               - SRKW conservation context mandatory
                               - No precise sensitive species coordinates
                               - Sacred species: nation-specific, never ceremonial
                               - Barred Owl: complex conservation accuracy
\end{asciibox}

\subsection{Safety Warden — Additional Gates (Part B)}

\begin{center}
\rowcolors{2}{rowgray}{white}
\begin{tabularx}{\textwidth}{L{2cm}XL{2.5cm}}
\toprule
\rowcolor{blockred}
\tblhdr{Test ID} & \tblhdr{Verifies} & \tblhdr{Boundary} \\
\midrule
BSC-01 & SRKW releases include ESA status, population, threats, NOAA recovery plan & BLOCK \\
BSC-02 & No precise coordinates for sensitive species (nesting eagles, salmon beds) & ABSOLUTE \\
BSC-03 & Sacred species (Eagle, Raven, Wolf) have nation-specific attribution & BLOCK \\
BSC-04 & Lummi Nation's Qwe'lhol mechen relationship to orca cited where relevant & GATE \\
BSC-05 & Barred Owl invasive status and USFWS management accurately described & GATE \\
BSC-06 & SOUNDSCAPE\_COMPOSITION entries framed as compositions, not species records & BLOCK \\
BSC-07 & Ecological crisis framing present but not catastrophizing & ANNOTATE \\
BSC-08 & All bioacoustic claims sourced (NOAA, Cornell Lab, peer-reviewed) & BLOCK \\
\bottomrule
\end{tabularx}
\end{center}

\subsection{SRKW — Southern Resident Killer Whale Mandatory Context}

\begin{warnbox}
\textbf{\sffamily\textcolor{secondary}{BLOCK-Level Requirement: SRKW Conservation Context}}

Every release at degrees 270–277, 311, 340, 348, and 357 MUST include:
\begin{itemize}[leftmargin=1.5em,topsep=2pt]
  \item Current SRKW population estimate (as of latest NOAA count — approximately 73 individuals, 2025)
  \item ESA Endangered Species Act listing status (Endangered, listed 2005)
  \item Primary threats: Chinook salmon prey depletion, vessel noise, contaminants
  \item Reference to NOAA's SRKW Recovery Plan (2023 update)
  \item Lummi Nation's cultural relationship (Qwe'lhol mechen — relatives in the sea)
\end{itemize}
The Safety Warden BLOCKS any SRKW release missing these elements.
Sources: NOAA Fisheries NWFSC; Lummi Nation (lummi-nsn.gov)
\end{warnbox}

\newpage

% ─────────────────────────────────────────────────────────────────────────────
% PART C — DEEP MAPPING
% ─────────────────────────────────────────────────────────────────────────────
\partheader{C}{Deep Mapping: A×B Cross-Reference Matrix}{colorC}

\section{Part C Vision: The Spaces Between}

\metafield{Requires:}{Part A complete (degrees 0–359) AND Part B complete (degrees 0–359)}
\metafield{Produces:}{360 A×B synthesis docs · 360 DIY Art project specs · cross-domain College nodes · skill-creator links}

\vspace{0.5em}

Part C is the moment the two half-circles meet. For each of the 360 degrees, Part C
loads both the music release and the species release, runs a Resonance Detector to find
structural parallels, synthesizes a cross-domain College lesson, generates a DIY Art
project spec that a student can actually execute, and links to GSD skill-creator resources.

\subsection{Selected A×B Resonance Pairs (Illustrative)}

\begin{center}
\rowcolors{2}{rowgray}{white}
\begin{tabularx}{\textwidth}{C{1.2cm}L{3.5cm}L{3.5cm}X}
\toprule
\rowcolor{colorC}
\tblhdr{Deg} & \tblhdr{Part A (Music)} & \tblhdr{Part B (Species)} & \tblhdr{Resonance Type} \\
\midrule
0 & Quincy Jones (ii–V–I, call-response) & Great Blue Heron (silence → croak) & Structured silence; call-and-response \\
1 & Bill Frisell (chromaticism) & Pacific Wren (syrinx polyphony) & Two-voice counterpoint in single organism/instrument \\
45 & Damien Jurado (folk unison) & Pacific Chorus Frog (emergent synchrony) & Rhythmic entrainment without a conductor \\
108 & Pileated Woodpecker drumming & Pileated Woodpecker (same!) & Degree 108 IS the woodpecker — Part A \& B merge \\
116 & American Dipper (song over noise) & American Dipper (same!) & Degree 116 is the Dipper — counterpoint against water \\
191 & Hovercraft (drone/noise) & Humpback Whale song (8-hour spiral) & Drone as temporal architecture \\
211 & Gray Wolf pack howl (A) & Gray Wolf pack howl (B) & Degrees 210–211 are the Wolf — Part A \& B merge \\
276 & Bikini Kill (matriarchal riot grrrl) & Orca female-led hunt sequence & Matriarchal coordination as acoustic structure \\
311 & Nirvana (community convergence) & Orca greeting ceremony (full pod) & Mass convergence call; collective identity assertion \\
332 & Botch (prime-number meters) & Woodpecker+River+Frog polyrhythm & Prime-number meters are a natural phenomenon \\
344 & Botch (peak mathcore) & Multi-species polyrhythm apex & Mathematical complexity at maximum biological intensity \\
359 & Unwound (post-hardcore silence) & Puget Sound (one breath) & Silence as the most compositionally charged moment \\
\bottomrule
\end{tabularx}
\end{center}

\subsection{Part C Component Architecture}

\begin{asciibox}
ab-matrix-deep-mapping/
|
+-- C00-ab-matrix-types --- ABPair, ResonanceNode, SynthesisDoc, DIYProject schemas
+-- C01-alignment-engine -- Load Part A + B releases -> ABPair[N] (Haiku)
+-- C02-resonance-detector - Find 2-5 structural parallels (Opus / HAWK)
|                            Types: Acoustic, Mathematical, Cultural, Energetic, Ecological
+-- C03-college-synthesizer - Produce cross-domain College lesson per degree (Sonnet / CEDAR)
|                             New dept: .college/cross-domain/
+-- C04-diy-art-generator --- Executable DIY art project spec per degree (Sonnet / WREN)
|                             Types: field-recording, instrument-build, graphic-score,
|                                    visual-art, collaborative-performance
+-- C05-skill-creator-linker - Link to GSD skill packs per degree (Haiku / OSPREY)
|                              foxfire-heritage, learn-kung-fu, deep-audio-analyzer, etc.
+-- C06-synthesis-warden ---- Gate: resonance accuracy + attribution (Opus / WARDEN)
                              BLOCK: reductive connection / romantic projection
                              GATE: Indigenous connection without specific citation
\end{asciibox}

\subsection{DIY Art Project Types (Per Degree)}

\begin{center}
\rowcolors{2}{rowgray}{white}
\begin{tabularx}{\textwidth}{L{3.5cm}XL{2.5cm}}
\toprule
\rowcolor{colorC}
\tblhdr{Project Type} & \tblhdr{Description} & \tblhdr{Example Degree} \\
\midrule
Field Recording & Go to specific PNW location; record species; edit to musical structure & Degree 45 (frog chorus) \\
Instrument Building & Build instrument that mimics specific acoustic property & Degree 109 (Pileated drum resonator) \\
Graphic Score & Create graphic notation of species call pattern & Degree 1 (Pacific Wren cascade) \\
Visual Art Prompt & Artwork expressing A×B resonance & Degree 276 (orca+Bikini Kill) \\
Collaborative Performance & Instructions for human + field recording performance & Degree 311 (Nirvana + orca) \\
Mathematics Art & Visualize the shared integer/frequency structure & Degree 332 (polyrhythm) \\
\bottomrule
\end{tabularx}
\end{center}

\subsection{New College Departments (Part C)}

\begin{center}
\rowcolors{2}{rowgray}{white}
\begin{tabularx}{\textwidth}{L{6cm}XC{2.5cm}}
\toprule
\rowcolor{colorC}
\tblhdr{Department Path} & \tblhdr{Content} & \tblhdr{Seeded By} \\
\midrule
\texttt{.college/cross-domain/} & Cross-disciplinary lesson plans (A×B resonances) & Part C \\
\texttt{.college/cross-domain/music-bioacoustics/} & Where music theory = bioacoustics & Part C \\
\texttt{.college/cross-domain/matriarchal-coordination/} & Leadership as acoustic structure & Degree 276 \\
\texttt{.college/cross-domain/prime-number-rhythm/} & Mathematics in natural percussion & Degrees 332, 344 \\
\texttt{.college/diy-art/field-recording/} & How to record PNW soundscapes & Part B→C |
\texttt{.college/diy-art/instrument-building/} & Build instruments from natural resonance & Part C \\
\texttt{.college/diy-art/graphic-score/} & Visual notation of biological sound & Part C \\
\bottomrule
\end{tabularx}
\end{center}

\newpage

% ─────────────────────────────────────────────────────────────────────────────
% PART D — FRACTAL REFINEMENT
% ─────────────────────────────────────────────────────────────────────────────
\partheader{D}{Fractal Refinement: Iterative Hi-Fi Quality Passes}{colorD}

\section{Part D Vision: Depth Without Bloat}

\metafield{Requires:}{Parts A, B, C complete}
\metafield{Principle:}{"A refinement pass must make the release more useful, not longer."}

\vspace{0.5em}

Part D is the fractal pass. A fractal is self-similar at every scale — the same patterns
that organize the 360° circle also organize each individual release. Part D looks inside
each release and finds the inner circle: the cross-references implied but not yet stated,
the mathematics bridges present but thin, the College nodes promoted when they have enough
examples, the DIY projects made executable when they were previously abstract.

\subsection{Refinement Pass Types}

\begin{center}
\rowcolors{2}{rowgray}{white}
\begin{tabularx}{\textwidth}{L{3.5cm}XL{2cm}L{3cm}}
\toprule
\rowcolor{colorD}
\tblhdr{Pass Type} & \tblhdr{Trigger} & \tblhdr{Model} & \tblhdr{What It Does} \\
\midrule
THEORY-DEEPEN & Concept in 5+ releases & Opus & Promote to full College lesson \\
CROSS-LINK & Release has $<$5 cross-refs & Sonnet & Add genealogy neighbor links \\
MATH-EXPLICIT & Math bridge thin & Opus & Add equations from \textit{The Space Between} \\
DIY-ENRICH & No materials list & Sonnet & Add materials, tools, sources \\
ATTRIBUTION-AUDIT & ANNOTATE in Pass 1 & Sonnet & Escalate ANNOTATE to fully resolved \\
COLLEGE-PROMOTE & Node has 8+ examples & Opus & Write full lesson module \\
\bottomrule
\end{tabularx}
\end{center}

\subsection{The Refinement Warden's Prime Directive}

\begin{amberbox}
\textbf{\sffamily\textcolor{colorC}{One Rule: Does This Pass Make the Release More Useful?}}

Before approving any refinement pass, the Warden asks: ``Would a student encountering this
release for the first time find it more useful after this pass?''

\medskip
\textbf{PASS:} The answer is yes, demonstrably.\\
\textbf{BLOCK:} The answer is ``more complete but harder to navigate.''\\
\textbf{BLOCK:} The answer is ``more equations but no pedagogical context.''\\
\textbf{BLOCK:} The answer is ``more cross-links but no explanation of why they connect.''

\medskip
Quality means clarity, precision, and utility. Not length. Not thoroughness. Usefulness.
\end{amberbox}

\subsection{Part D Component Architecture}

\begin{asciibox}
fractal-refinement-engine/
|
+-- D00-refinement-types --- RefinementPass, QualityScore, PromotionCandidate schemas
+-- D01-quality-auditor ---- Score all 720 releases on 10-point rubric (Sonnet)
|                            Dimensions: theory-depth, math-bridge, cross-links,
|                                        diy-executability, attribution-completeness
+-- D02-fractal-expander ---- Deepen qualifying releases (Opus)
|                             Add Space Between equations, theory genealogy depth
+-- D03-cross-linker -------- Add intra-release cross-refs; build genealogy graph (Sonnet)
+-- D04-college-promoter ---- Promote 8+-example nodes to full lessons (Opus)
+-- D05-refinement-warden --- Gate: improvement quality (Opus)
                              BLOCK: longer-but-less-useful
                              BLOCK: equations without pedagogy
\end{asciibox}

\newpage

% ─────────────────────────────────────────────────────────────────────────────
% PART E — LESSONS APPLIED
% ─────────────────────────────────────────────────────────────────────────────
\partheader{E}{Lessons Applied: What We Do Better Next Rotation}{colorE}

\section{Part E Vision: The Meta-Retrospective}

\metafield{Requires:}{Parts A, B, C, D all complete}
\metafield{Produces:}{Meta-retrospective · Cycle 2 specification · Complete Edition PDF}

\vspace{0.5em}

Every rotation teaches the next one. Part E is not a summary. It is an architectural audit
with the specificity of someone who just ran 77.6 million tokens of work across 1,800 artifacts
and wants to make absolutely sure the next cycle doesn't repeat the same mistakes.

\subsection{Key Questions Part E Answers}

\begin{center}
\rowcolors{2}{rowgray}{white}
\begin{tabularx}{\textwidth}{L{3cm}X}
\toprule
\rowcolor{colorE}
\tblhdr{Domain} & \tblhdr{Questions} \\
\midrule
Theory & Which music theory concepts proved most teachable through A×B pairing? \\
Bioacoustics & Which species produced the richest content? Which were scientifically thin? \\
College & Which departments grew beyond expectations? What needs restructuring? \\
Pipeline & Where did token budget deviate most? Which waves bottlenecked? \\
Safety & Which BLOCK-level findings recurred? What structural Safety Warden improvements? \\
DIY Art & Which project specs were executable? Which were too abstract? \\
Culture & Where did OCAP® compliance require the most careful handling? \\
\bottomrule
\end{tabularx}
\end{center}

\subsection{Part E Outputs}

\begin{center}
\rowcolors{2}{rowgray}{white}
\begin{tabularx}{\textwidth}{L{5cm}XL{3.5cm}}
\toprule
\rowcolor{colorE}
\tblhdr{Output} & \tblhdr{Description} & \tblhdr{Audience} \\
\midrule
\texttt{meta-retrospective.md} & Patterns across all 720 retrospectives & GSD architect \\
\texttt{lessons-learned-taxonomy.md} & Categorized lessons: theory/bio/college/pipeline/safety & All future cycles \\
\texttt{cycle-2-specification.md} & Full Vision doc for "Second Rotation" & GSD orchestrator \\
\texttt{ecosystem-updates.md} & Specific improvements to skill-creator + College & Claude Code \\
\texttt{sound-of-puget-sound-complete.pdf} & Complete Edition — all 5 parts, publication quality & Public / College \\
\bottomrule
\end{tabularx}
\end{center}

\newpage

% ─────────────────────────────────────────────────────────────────────────────
% MISSION PACKAGE — ALL PARTS
% ─────────────────────────────────────────────────────────────────────────────
\stageheader{STAGE 3 \quad COMPLETE MISSION PACKAGE — ALL FIVE PARTS}{primary}

\section{Complete Component Matrix}

\begin{center}
\rowcolors{2}{rowgray}{white}
\begin{tabularx}{\textwidth}{C{0.5cm}L{4.5cm}C{1.5cm}C{1.8cm}C{2.2cm}X}
\toprule
\rowcolor{primary}
\tblhdr{\#} & \tblhdr{Component} & \tblhdr{Part} & \tblhdr{Wave} & \tblhdr{Model} & \tblhdr{Agent} \\
\midrule
A0–A8 & 9 components (Part A) & A & 0–3d & Haiku/Sonnet/Opus & RAVEN/SALMON/CEDAR/ORCA/HAWK/WARDEN/HERON/OWL/D-FIR \\
B0–B8 & 9 components (Part B) & B & 0–3d & Haiku/Sonnet/Opus & RAVEN/SALMON/EAGLE/ORCA/HAWK/WARDEN/HERON/OWL/CEDAR \\
C0–C6 & 7 components (Part C) & C & 0–3 & Haiku/Sonnet/Opus & RAVEN/SALMON/HAWK/CEDAR/WREN/OSPREY/WARDEN \\
D0–D5 & 6 components (Part D) & D & 0–3 & Sonnet/Opus & RAVEN/AUDITOR/HAWK/CEDAR/PROMOTER/WARDEN \\
E0–E5 & 6 components (Part E) & E & 0–3 & Sonnet/Opus & RAVEN/SALMON/HAWK/CEDAR/OSPREY/WARDEN \\
\midrule
\rowcolor{lightprimary}
\textbf{37} & \textbf{Total Components} & \textbf{A–E} & & & \\
\bottomrule
\end{tabularx}
\end{center}

\section{Complete Test Matrix}

\begin{center}
\rowcolors{2}{rowgray}{white}
\begin{tabularx}{\textwidth}{L{2cm}C{2cm}C{2.5cm}C{2.5cm}X}
\toprule
\rowcolor{primary}
\tblhdr{Part} & \tblhdr{Total Tests} & \tblhdr{Safety-Critical} & \tblhdr{Failure Action} & \tblhdr{Key Safety Gates} \\
\midrule
A & 148 & 22 & BLOCK & Attribution, OCAP®, compilation \\
B & 152 & 24 & BLOCK & SRKW context, location, sacred species \\
C & 120 & 18 & BLOCK & Resonance accuracy, Indigenous attribution \\
D & 90 & 12 & BLOCK & Quality not length, pedagogy not equations \\
E & 60 & 15 & BLOCK & Meta-accuracy, Cycle 2 completeness \\
\midrule
\rowcolor{lightprimary}
\textbf{Total} & \textbf{570} & \textbf{91} & & \\
\bottomrule
\end{tabularx}
\end{center}

\section{Complete Execution Summary}

\begin{center}
\rowcolors{2}{rowgray}{white}
\begin{tabularx}{\textwidth}{L{6cm}X}
\toprule
\rowcolor{primary}
\tblhdr{Metric} & \tblhdr{Value} \\
\midrule
Total releases & 1,800 artifacts (360+360+360+720+5) \\
Total components & 37 across Parts A–E \\
Total tests & 570 (91 safety-critical mandatory-pass) \\
Activation profile & Fleet+ (Parts A/B) · Fleet (Parts C/D/E) \\
Model split (overall) & $\approx$35\% Opus / $\approx$55\% Sonnet / $\approx$10\% Haiku \\
Inputs & 2 CSVs: 360 artists + 360 species \\
New College departments & 7 new departments (bioacoustics, diy-art, cross-domain, etc.) \\
College nodes (est.) & $\approx$2,500 nodes created or enriched \\
DIY Art projects & 360 executable project specs \\
Token budget (total) & $\approx$75.5M across all 5 parts \\
Context windows & $\approx$2,800 total \\
Estimated sessions & $\approx$700 (at 4 windows/session) \\
A+B parallel execution & Yes — Parts A and B run concurrently \\
CAPCOM escalation & Per-release BLOCK findings (human review required) \\
Completion signal & Part E: \texttt{CYCLE-COMPLETE + CYCLE-2-READY} \\
Output root & \texttt{releases/sound-of-puget-sound/} \\
\bottomrule
\end{tabularx}
\end{center}

\section{Execution Order}

\begin{asciibox}
Phase 1 (CONCURRENT):
  [Part A engine] -----> Degrees 0-359 (seattle-360)
  [Part B engine] -----> Degrees 0-359 (pnw-360)
  Both pipelines run simultaneously. No shared state.
  Completion: Both reach degree 359 and emit CYCLE-COMPLETE.

Phase 2 (After both A + B at 359):
  [Part C engine] -----> Degrees 0-359 (ab-matrix-deep-mapping)
  Requires: releases/seattle-360/ AND releases/pnw-360/ fully populated

Phase 3 (After Part C at degree 359):
  [Part D engine] -----> Quality passes on releases/seattle-360/ + releases/pnw-360/
  Runs 720 refinement passes (not a 360-degree loop -- a quality audit loop)

Phase 4 (After Part D quality audit complete):
  [Part E engine] -----> Meta-retrospective + Cycle 2 specification
  Produces: sound-of-puget-sound-complete.pdf (the final artifact)

Total wall time estimate: ~700 sessions across ~2,800 context windows
Parts A+B together: ~360 sessions (concurrent)
Part C: ~180 sessions
Part D: ~120 sessions
Part E: ~40 sessions
\end{asciibox}

\vspace{1.5em}

\begin{quotebox}
\textbf{\sffamily\textcolor{primary}{The Through-Line:}}

\medskip
\textit{``The Pacific Wren holds the record for the most complex song relative to body size of
any bird in North America. It has been singing this way, in the old-growth understory of the
Pacific Northwest, since before humans arrived. It will continue doing so, if we leave enough
old-growth standing.}

\medskip
\textit{This engine is built to trace the thread that connects the Wren's song to Quincy Jones's
chord voicing, that connects the orca's greeting ceremony to Bikini Kill's matriarchal coordination,
that connects Botch's prime-number meters to the stonefly hatch's wing-buzz polyrhythm.}

\medskip
\textit{The 360° rotation is complete when a student can hear the Pacific Wren, understand the
interval ratios in its song, trace those ratios to \textit{The Space Between}'s unit circle, find the
parallel in Quincy Jones's chord voicings, and feel the full circle close in their understanding.}

\medskip
\textit{At that moment, the Pacific Northwest has given back what it always had: the sound
of everything, teaching everything, to anyone willing to listen.}

\medskip
\textit{One degree at a time. It's all about the journey.''}

\medskip
\textcolor{subtlegray}{\sffamily\small The Sound of Puget Sound · A–B–C–D–E Complete Release Matrix · GSD Ecosystem · 2026}
\end{quotebox}

\end{document}
